\documentclass[12pt,a4paper]{article}

\usepackage[T1]{fontenc}
\usepackage[utf8]{inputenc}
\usepackage[ngerman]{babel}
\usepackage{lmodern}
\usepackage{microtype}

\usepackage{amsmath,amssymb,amsthm,mathtools}
\usepackage[a4paper,margin=2.4cm]{geometry}
\usepackage{booktabs}
\usepackage{xcolor}
\usepackage[hidelinks,unicode]{hyperref}
\usepackage[most]{tcolorbox}
\usepackage{enumitem}

\definecolor{bewiesengruen}{RGB}{30,100,50}
\definecolor{warnorange}{RGB}{200,100,10}
\definecolor{lueckenrot}{RGB}{180,40,40}
\definecolor{kernblau}{RGB}{20,60,130}

\tcbset{
  bewiesen/.style={colback=green!7,colframe=bewiesengruen,fonttitle=\bfseries},
  heuristik/.style={colback=orange!10,colframe=warnorange,fonttitle=\bfseries},
  offen/.style={colback=red!8,colframe=lueckenrot,fonttitle=\bfseries},
  kern/.style={colback=blue!5,colframe=kernblau,fonttitle=\bfseries},
}

\newtheorem{satz}{Satz}[section]
\newtheorem{lemma}[satz]{Lemma}
\newtheorem{definition}[satz]{Definition}
\newtheorem{bemerkung}[satz]{Bemerkung}
\newtheorem{vermutung}[satz]{Vermutung}

\newcommand{\vii}{\nu_2}
\newcommand{\N}{\mathbb{N}}
\newcommand{\Z}{\mathbb{Z}}
\newcommand{\Ztwo}{\mathbb{Z}_2}
\newcommand{\Ztwost}{\mathbb{Z}_2^{\ast}}
\newcommand{\U}{\mathcal{U}}
\newcommand{\disttwo}{\operatorname{dist}_2}
\newcommand{\Einf}{E^{\infty}}
\newcommand{\Ediag}{E^{\mathrm{diag}}}
\newcommand{\Etail}{E^{\mathrm{tail}}}
\newcommand{\Ecap}{E^{\cap}}

\title{%
  Äquivalenzkette Collatz $\leftrightarrow$ $E^{\infty}$ $\leftrightarrow$\\
  EABC-Grammatik
}
\author{Thomas Hoffbauer}
\date{17.\ Juni 2026}

\begin{document}
\maketitle

\begin{abstract}
Exakte Definitionen für odd-to-odd-Collatz, $E_{N,K}$, $\Ediag$ und
$\Einf$; Äquivalenz Collatz $\Leftrightarrow$ $\Einf\cap\N=\emptyset$
und Formulierung~B; Nicht-Äquivalenz von $\Ediag$ (Beispiel $27$);
Skizze Satz~L und EABC-Grammatik. \textbf{Collatz nicht bewiesen.}
\end{abstract}

% =================================================
\section{Einleitung}
\label{sec:einleitung}
% =================================================

Die klassische Collatz-Vermutung verlangt punktweise Konvergenz:
jede Bahn erreicht den Zyklus $(1,2,4)$.
Bewiesen sind dagegen nur \emph{mittlere} oder \emph{fast-sichere}
Aussagen (z.\,B.\ Tao 2019: logarithmische Dichte~$1$ für beschränkte
Bahnen). Der Engpass ist die Brücke
\[
  \text{Mittel / Markov / Drift}\quad\nRightarrow\quad\forall n\in\N.
\]

In der 2-adischen Geometrisierung wird eine abgeschlossene Ausnahmemenge
$E\subset\Ztwo$ als Limes endlicher Approximationen konstruiert.
Die naheliegende Wahl
\[
  E\;:=\;\Ediag\;:=\;\overline{\bigcup_{N\in\N} E_{N,N}}^{\;\Ztwo}
\]
ist \textbf{nicht} äquivalent zur Collatz-Vermutung.
Beispiel $n=27$: Die odd-to-odd-Stoppzeit $\tau(27)$ liegt bei
$\approx 41$ Schritten (deutlich über $27$).
Für $N=K=27$ gilt daher $27\in E_{27,27}$ und folglich
$27\in\Ediag\cap\N$ --- obwohl Collatz für $27$ wahr ist
(die Bahn erreicht $1$).
Der Diagonal-Limes misst \emph{endliche Beobachtungsfenster},
nicht ewige Nicht-Konvergenz.

% =================================================
\section{Definitionen}
\label{sec:def}
% =================================================

Sei $\Ztwo$ der Ring der 2-adischen ganzen Zahlen.
Jedes ungerade $n\in\N$ identifizieren wir mit seinem Bild in $\Ztwo$.

\begin{definition}[Odd-to-odd-Abbildung und Stoppzeit]
\label{def:U}
Für ungerades $x\in\Ztwo$ mit $3x+1\neq 0$ setzen wir
\[
  \U(x):=\frac{3x+1}{2^{\vii(3x+1)}}\in\Ztwo.
\]
Auf $\N$ liefert $\U$ die übliche odd-to-odd-Collatz-Kette.
Die \emph{Stoppzeit} (falls sie existiert) ist
\[
  \tau(n):=\min\{K\geq 1:\ \U^K(n)=1\}\in\N\cup\{\infty\}.
\]
\end{definition}

\begin{definition}[Endliche Ausnahmeapproximationen]
\label{def:ENK}
Für $N,K\in\N$ sei
\[
  E_{N,K}\ :=\ \bigl\{x\in\Ztwo:\ \exists\,\text{ungerades }n\leq N,\ (n)_{\Ztwo}=x,\
  \ \forall k\leq K:\ \U^k(n)\neq 1\bigr\}.
\]
Monotonie: $N_1\leq N_2\Rightarrow E_{N_1,K}\subseteq E_{N_2,K}$;
$K_1\leq K_2\Rightarrow E_{N,K_2}\subseteq E_{N,K_1}$.
Der \emph{Diagonal-Akkumulator} und sein Abschluss sind
\[
  E^{\mathrm{acc}}\ :=\ \bigcup_{N\in\N} E_{N,N},
  \qquad
  \Ediag\ :=\ \overline{E^{\mathrm{acc}}}^{\;\Ztwo}.
\]
\end{definition}

\begin{definition}[$E^{\infty}$ und Hilfsmengen]
\label{def:Einfty}
\[
  \Einf\ :=\ \bigl\{n\in\N:\ n\ \text{ungerade},\ \forall K\in\N:\ \U^K(n)\neq 1\bigr\}
  \ =\ \{\,n\ \text{ungerade}:\ \tau(n)=\infty\,\}.
\]
Kurz erwähnt:
\begin{itemize}[nosep]
  \item $\Ecap:=\bigcap_{N} E_{N,N}$; für $n\in\N$ gilt
    $n\in\Ecap\cap\N\Leftrightarrow\tau(n)=\infty$, also
    $\Ecap\cap\N=\Einf$.
  \item $\Etail\subset\Ztwo$: Punkte, deren Bahn den trivialen Bereich
    $A_{\mathrm{triv}}:=\{x:\ |x-1|_2<1\}$ in keinem endlichen
    Präfix trifft (2-adische Tail-Ausnahme).
\end{itemize}
\end{definition}

\begin{definition}[2-adische Distanz]
\label{def:dist2}
Für $x,y\in\Ztwo$, $x\neq y$: $|x-y|_2:=2^{-\vii(x-y)}$.
Für $n\in\N$, $E\subset\Ztwo$:
$\disttwo(n,E):=\inf_{e\in E}|n-e|_2$,
mit $\disttwo(n,\emptyset)=1$.
\end{definition}

\begin{bemerkung}[Warum $\disttwo(\cdot,1)$ scheitert]
\label{rem:dist1}
Die naive Uniformität $\disttwo(n,1)\geq c\cdot 2^{-\lfloor\log_2 n\rfloor}$
ist \textbf{widerlegt}: Für $n_0=2^{k+1}3^r-1$ bleibt
$\disttwo(n_0,1)$ oft bei $1/2$ konstant (LTE-Familie), während
$c\cdot 2^{-\lfloor\log_2 n\rfloor}\to 0$.
Formal: \texttt{dist\_to\_one\_not\_uniform\_bound} in
\texttt{collatz\_z2\_attraktor.lean}.
Der Attraktor ist $A_{\mathrm{triv}}$ bzw.\ eine echte Ausnahmemenge,
nicht der isolierte Punkt $1\in\Ztwo$.
\end{bemerkung}

% =================================================
\section{Äquivalenzkette}
\label{sec:aek}
% =================================================

\begin{satz}[Collatz $\Leftrightarrow$ $E^{\infty}$]
\label{sat:collatz-einf}
Sei \emph{Collatz (odd-to-odd)} die Aussage:
$\forall$ ungerade $n\in\N\ \exists K:\ \U^K(n)=1$.
Dann gilt
\[
  \text{Collatz (odd-to-odd)}\quad\Longleftrightarrow\quad
  \Einf\cap\N=\emptyset.
\]
\end{satz}

\begin{proof}
$\Rightarrow$: Ist Collatz wahr, existiert für jedes ungerade $n$ ein $K$
mit $\U^K(n)=1$, also $n\notin\Einf$.

$\Leftarrow$: Ist $\Einf\cap\N=\emptyset$, hat jedes ungerade $n$ endliche
Stoppzeit $\tau(n)<\infty$, d.\,h.\ $\U^{\tau(n)}(n)=1$.
\end{proof}

\begin{satz}[Collatz $\Leftrightarrow$ Formulierung~B]
\label{sat:formB}
Collatz (odd-to-odd) ist äquivalent zu
\[
  \forall n\in\N\ \exists K\in\N:\quad n\notin E_{n,K}.
\]
\end{satz}

\begin{proof}
Für ungerades $n$ gilt $n\in E_{n,K}$ genau dann, wenn
$\forall k\leq K:\ \U^k(n)\neq 1$, also genau wenn $\tau(n)>K$.
Daher $\exists K:\ n\notin E_{n,K}\Leftrightarrow\tau(n)<\infty$.
\end{proof}

\begin{satz}[Nicht-Äquivalenz: $\Ediag$ vs.\ Collatz]
\label{sat:nondiag}
Unter der Collatz-Vermutung gilt $\Ediag\cap\N\neq\emptyset$.
Insbesondere ist $\Ediag\cap\N=\emptyset$ \textbf{streng schwächer}
als Collatz und nicht äquivalent.
\end{satz}

\begin{proof}
Sei $n=27$. Dann $\tau(27)>27$, also $27\in E_{27,27}\subseteq E^{\mathrm{acc}}$
und damit $27\in\Ediag\cap\N$.
Unter Collatz ist dennoch $\tau(27)<\infty$, also $27\notin\Einf$.
\end{proof}

\begin{vermutung}[Zusammenhang $\Ediag$ und $\Einf$ in $\Ztwo$]
\label{verm:z2}
In $\Ztwo$ gilt $\Einf\subseteq\Ediag$ (jede natürliche Ausnahme liegt im
Diagonal-Limes). Die Umkehrung und die Relation zu $\Etail$ sind offen:
Ist $\Ediag\cap\N$ genau die Menge der $n$ mit $\tau(n)>n$ für
\emph{irgendwelche} endlichen Fenster, oder gibt es 2-adische Grenzpunkte
in $\Ediag\setminus\N$, die Collatz nicht reflektieren?
\end{vermutung}

\begin{tcolorbox}[kern,title={Formulierungen}]
Collatz-äquivalent: $\Einf\cap\N=\emptyset$;
$\forall n\ \exists K:\ n\notin E_{n,K}$.
Nicht äquivalent: $\Ediag\cap\N=\emptyset$ (Satz~\ref{sat:nondiag});
$\disttwo(n,1)\geq c\cdot 2^{-\log n}$ (Bem.~\ref{rem:dist1}).
\end{tcolorbox}

% =================================================
\section{Satz~L (Zwischenlemma)}
\label{sec:satzL}
% =================================================

\begin{vermutung}[Satz~L: präperiodischer EABC-Tail]
\label{verm:satzL}
Sei $w(n)=(w_k)_{k\geq 0}\in\{E,A,B,C\}^{\N}$ die EABC-Paritätsfolge
der odd-to-odd-Bahn von $n$ (Restklassen $\bmod 12$ der ungeraden
Zwischenwerte). Dann ist $w(n)$ \emph{präperiodisch} im Sinne:
\[
  \forall N_0\in\N\ \exists k\geq N_0,\ p\geq 1:\quad
  w_{k+j}=w_{k+j+p}\quad\text{für alle }j\geq 0
  \text{ bis zum ersten }K\text{ mit }\U^K(n)=1.
\]
\end{vermutung}

\begin{bemerkung}[Plausibilität]
Bewiesen: endliche C-Ketten, Verbot B$\to$B, EA-Zwang nach C.
Satz~L wäre asymptotische Wort-Regularität vor Erreichen von~$1$.
\textbf{Status: offen;} aus Satz~L folgt Collatz nicht ohne Uniformität.
\end{bemerkung}

% =================================================
\section{Grammatik und symbolische Dynamik}
\label{sec:grammatik}
% =================================================

\begin{definition}[EABC-Sprache]
\label{def:grammatik}
Alphabet $\mathcal{A}=\{E,A,B,C\}$ mit
$E\equiv 1$, $A\equiv 5$, $B\equiv 7$, $C\equiv 11\pmod{12}$.
Sei $\mathcal{L}\subset\mathcal{A}^*$ die Menge der \emph{realisierbaren}
endlichen Wörter, die als EABC-Folge einer odd-to-odd-Collatz-Bahn
vorkommen. Zulässigkeit (bewiesen bzw.\ strukturell):
\begin{enumerate}[label=(\roman*),nosep]
  \item kein Teilwort $BB$;
  \item C-Ketten endlich: nach $C^k$, $k\geq 1$, folgt kein weiteres $C$
    über die maximale C-Kettenlänge hinaus;
  \item EA-Zwang: nach jeder maximalen C-Kette tritt ein Block mit
    $\vii\geq 2$ auf (EA-Schritt).
\end{enumerate}
\end{definition}

\begin{satz}[Grammatik-Formulierung (konzeptionell)]
\label{sat:grammatik}
Collatz (odd-to-odd) ist äquivalent zur Aussage:
Jede unendliche Realisierung $w\in\mathcal{A}^{\N}$, die von einem
ungeraden $n\in\N$ erzeugt wird, enthält nur endlich viele Symbole
vor dem ersten Auftreten des Zyklus $E$-Tail zum Fixpunkt~$1$;
äquivalent: es gibt keine unendliche realisierbare Wortfolge in
$\mathcal{L}$, deren zugehörige Bahn $\U^K(n)\neq 1$ für alle $K$ erfüllt.
\end{satz}

\begin{bemerkung}
Die 2-adische Shift-Konjugation
$\Phi\circ S\circ\Phi^{-1}=T$ auf $\Ztwo$
~\cite{BernsteinLagarias1996,Bernstein1994} und die Paritäts-Isometrie
$Q$~\cite{Rozier2019,Lagarias1985} liefern den topologischen Rahmen:
Collatz wird zur Frage nach zulässigen unendlichen Paritätswörtern
in einer Untersprache von $\{0,1\}^{\N}$, gefiltert durch EABC-Regeln.
\end{bemerkung}

% =================================================
\section{Epilog}
\label{sec:epilog}
% =================================================

Die 2-adische Geometrisierung verbindet die Äquivalenzkette mit
EABC-Struktur, LTE-Kompensation, Markov-Mischung auf $\{E,A,B,C\}$
und dem Grenzobjekt $\Einf\subset\Ediag\subset\Ztwo$.
Bewiesen sind lokale Grammatikregeln und Dichteaussagen;
\textbf{nicht bewiesen} ist die punktweise Brücke
\[
  \text{Struktur auf Blöcken}\quad\Longrightarrow\quad
  \Einf\cap\N=\emptyset.
\]

\begin{tcolorbox}[offen,title={Verbleibende Lücke}]
Zwischen Formulierung~B ($\forall n\ \exists K:\ n\notin E_{n,K}$)
und dem Diagonal-Limes $\Ediag$ liegt genau die Uniformitätsfrage:
\emph{wie} wächst $K(n)$ mit $n$, und gibt es 2-adische Grenzpunkte,
die endliche Fenster täuschen, ohne in $\Einf$ zu liegen?
Satz~L wäre ein Zwischenschritt; sein Beweis oder seine Widerlegung
würde die Grammatik- und $E^{\infty}$-Formulierung schärfen.
\end{tcolorbox}

% =================================================
\begin{thebibliography}{99}
% =================================================

\bibitem{BernsteinLagarias1996}
D.~J.~Bernstein und J.~C.~Lagarias.
The $3x+1$ conjugacy map.
\textit{Canad.\ J.\ Math.} \textbf{48} (1996), 1154--1169.

\bibitem{Bernstein1994}
D.~J.~Bernstein.
A non-iterative $2$-adic statement of the $3x+1$ conjecture.
\textit{Proc.\ Amer.\ Math.\ Soc.} \textbf{121} (1994), 405--408.

\bibitem{Rozier2019}
O.~Rozier.
Parity sequences of the $3x+1$ map on the $2$-adic integers and Euclidean embedding.
\textit{Integers} \textbf{19} (2019), A8.
arXiv:1805.00133.

\bibitem{Lagarias1985}
J.~C.~Lagarias.
The $3x+1$ problem and its generalizations.
\textit{Amer.\ Math.\ Monthly} \textbf{92} (1985), 3--23.

\end{thebibliography}

\end{document}
